\SourcePage{294}
\SGAchapter{VII}{若干经典概形的上同调与陈类的上同调理论}{Cohomologie de quelques schémas classiques et théorie cohomologique des classes de Chern}
\begin{center}J. P. Jouanolou\end{center}

全文固定一个整数 $\nu>0$，并置 $A=\mathbf Z/\nu\mathbf Z$；所考虑的所有\Term{scheme}{概形}的\Term{residue-characteristic}{剩余特征}均与 $\nu$ 互素。若 $X$ 是这样的\Term{scheme}{概形}，以 $\mu_X$ 表示 $\nu$ 次\Term{root-of-unity}{单位根}\Term{sheaf}{层}，以 $A^\bullet(X)$ 表示\Term{cohomology-ring}{上同调环}
\[
 \bigoplus_i H^{2i}(X,\mu_X^{\otimes i}),
\]
其中\Term{multiplication}{乘法}为\Term{cup-product}{杯积}。

\SGAsection{1}{向量丛}
设 $S$ 是一个\Term{scheme}{概形}，$E$ 是\Term{locally-free}{局部自由}的 $\mathscr O_S$-\Term{module}{模}（默认为\Term{finite-type}{有限型}）。置 $X=\mathbf V(E)$、$U=\mathbf V(E)^*$（即\Term{zero-section}{零截面}的补集），
\SourcePage{295}
并以 $p:X\to S$ 和 $q:U\to S$ 表示\Term{canonical-projection}{典范投影}。

\begin{SGAresult}{命题}{1.1}
设 $L$ 是一个 $A$-\Term{module}{模}。

(i) a) 对每个整数 $j>0$，有 $R^jp_*(p^*L)=0$。

b) \Term{adjunction-morphism}{伴随态射} $L\to p_*p^*(L)$ 是\Term{isomorphism}{同构}。

(ii) 假设 $E$ 的\Term{rank}{秩}恒为 $r$。

a) 当 $i\ne2r$ 时，$R^ip_!(p^*L)=0$。

b) \Term{trace-morphism}{迹态射}（SGA 4 第 XVIII 讲）
\[
 R^{2r}p_!(p^*L)\longrightarrow L\otimes\mu_S^{\otimes-r}
\]
是\Term{isomorphism}{同构}。
\end{SGAresult}
\begin{proof}
(i) 对 Zariski \Term{topology}{拓扑}作\Term{localization}{局部化}后，立即由 SGA 4 第 XV 讲 2.2 得出。(ii) 的内容基本上已包含在 SGA 4 第 XVIII 讲中。
\end{proof}

\begin{SGAresult}{推论}{1.2}
设 $L$ 是一个 $A_S$-\Term{module}{模}。对每个整数 $i$，\Term{inverse-image}{逆像}\Term{morphism}{态射}
\[
 H^i(S,L)\longrightarrow H^i(X,p^*L)
\]
是\Term{isomorphism}{同构}。
\end{SGAresult}
\begin{proof}
由 1.1，\Term{adjunction-morphism}{伴随态射} $L\to\mathbf Rp_*(p^*L)$ 是\Term{isomorphism}{同构}。又有 $H^*(X,p^*L)=H^*(S,\mathbf Rp_*p^*(L))$，故结论成立。
\end{proof}

\begin{SGAresult}{命题}{1.3}
假设 $E$ 的\Term{rank}{秩}恒为 $r$。对每个 $A_S$-\Term{module}{模} $L$，下列性质成立。

(i) a) 当 $i\ne0,2r-1$ 时，$R^iq_*(q^*L)=0$。

b) \Term{adjunction-morphism}{伴随态射} $L\to q_*q^*(L)$ 是\Term{isomorphism}{同构}。

c) 以 $Y$ 表示 $U$ 在 $X$ 中的补集，则\Term{morphism}{态射}
\[
 R^{2r-1}q_*(q^*L)\longrightarrow L\otimes\mu_S^{\otimes-r},
\]
\SourcePage{296}
即“\Term{boundary-morphism}{边界态射}”
\[
 R^{2r-1}q_*(q^*L)\longrightarrow R_Y^{2r}p_*(p^*L)\tag{1}
\]
与\Term{purity-isomorphism}{纯性同构}（SGA 4 第 XVIII 讲）
\[
 R_Y^{2r}p_*(p^*L)\longrightarrow L\otimes\mu_S^{\otimes-r}
\]
的复合，是\Term{isomorphism}{同构}。

(ii) a) 当 $i\notin\{1,2r\}$ 时，$R^iq_!(q^*L)=0$。

b) 从双向无限的\Term{exact-sequence}{正合列}（$f:Y\to S$ 为\Term{canonical-projection}{典范投影}）
\[
 \cdots\longrightarrow R^iq_!(q^*L)\longrightarrow R^ip_!(p^*L)
 \longrightarrow R^if_!(f^*L)\longrightarrow R^{i+1}q_!(q^*L)\longrightarrow\cdots
\]
得到的\Term{boundary-morphism}{边界态射} $L\to R^1q_!(q^*L)$ 是\Term{isomorphism}{同构}。

c) \Term{trace-morphism}{迹态射}
\[
 R^{2r}q_!(q^*L)\longrightarrow L\otimes\mu_S^{\otimes-r}
\]
是\Term{isomorphism}{同构}。
\end{SGAresult}
\begin{proof}
设 $\alpha:U\to X$ 和 $\beta:Y\to X$ 为\Term{canonical-immersion}{典范浸入}。置 $F=p^*(L)$，已知（SGA 4 第 XVI 讲 3.7）
\[
 F\xrightarrow{\sim}\alpha_*\alpha^*(F),\qquad
 (R^j\alpha_*)\alpha^*(F)=0\quad(j\ne0,2r-1).
\]
另一方面，当 $j\ge1$ 时，\Term{sheaf}{层} $(R^j\alpha_*)(\alpha^*F)$ 的\Term{support}{支集}包含于 $Y$，所以对于\Term{functor}{函子} $p_*$ 是\Term{acyclic}{无上同调的}。利用 Leray \Term{spectral-sequence}{谱序列}
\[
 R^ip_*(R^j\alpha_*)(q^*L)\Longrightarrow R^nq_*(q^*L),
\]
由这些观察即得 (i) a)、b)。断言 (i) c) 则由 1.1 (i) 以及双向无限的\Term{exact-sequence}{正合列}
\[
 \cdots\longrightarrow R_Y^ip_*(p^*L)\longrightarrow R^ip_*(p^*L)
 \longrightarrow R^iq_*(q^*L)\longrightarrow R_Y^{i+1}p_*(p^*L)\longrightarrow\cdots
\]
得到。\footnote{\Term{sheaf}{层} $R_Y^{2r}p_*(p^*L)$ 是与\Term{presheaf}{预层} $U\mapsto H^{2r}_{Y\times_SU}(X\times_SU,p_U^*L)$ 相伴的\Term{sheaf}{层}。}

\SourcePage{297}
再证明 (ii)。c) 来自 $q$ 是\Term{smooth}{光滑}\Term{morphism}{态射}，且其\Term{geometric-fiber}{几何纤维}\Term{connected}{连通}、\Term{dimension}{维数}为 $r$。b) 中的\Term{exact-sequence}{正合列}由对\Term{canonical}{典范}\Term{exact-sequence}{正合列}
\[
 0\longrightarrow\alpha_!(q^*L)\longrightarrow p^*L
 \longrightarrow\beta_!(\beta^*p^*L)\longrightarrow0
\]
应用\Term{functor}{函子} $\mathbf Rp_!$ 得到。利用 1.1 (ii)，由所得双向无限的\Term{exact-sequence}{正合列}容易\Term{pushout}{推出} (ii) a)、b)。
\end{proof}

\textbf{注记 1.4。} 当 $S=\operatorname{Spec}(\mathbf C)$ 时，1.3 表明 $\mathbf C^r$ 去掉原点所得\Term{space}{空间}，与实\Term{dimension}{维数}为 $2r-1$ 的球面有相同的\Term{cohomology}{上同调}；该球面是 $\mathbf C^r\setminus0$ 的\Term{deformation-retract}{形变收缩核}。

\begin{SGAresult}{推论（Gysin \Term{exact-sequence}{正合列}）}{1.5}
Leray \Term{spectral-sequence}{谱序列}
\[
 H^i(S,R^jq_*(q^*L))\Longrightarrow H^{i+j}(U,q^*L)
\]
给出双向无限的\Term{exact-sequence}{正合列}（称为 Gysin \Term{exact-sequence}{正合列}）
\[
 \cdots\longrightarrow H^i(S,L)\longrightarrow H^i(U,q^*L)
 \longrightarrow H^{i-2r+1}(S,L\otimes\mu_S^{\otimes-r})
 \longrightarrow H^{i+1}(S,L)\longrightarrow\cdots.
\]
\end{SGAresult}
\begin{proof}两行\Term{lemma}{引理}。\end{proof}

\SGAsection{2}{射影概形}
以下的叙述方式基于 L. Illusie 的建议，所用方法与 P. Deligne 为证明某些\Term{spectral-sequence}{谱序列}\Term{degeneration}{退化}而采用的方法 [6] 相近。

\textbf{2.1。} 设 $X,Y$ 是两个\Term{scheme}{概形}，$f:X\to Y$ 是一个\Term{morphism}{态射}。我们将定义 $D(Y)$ 中的\Term{arrow}{箭头}
\[
 \mathbf Rf_*(A_X)\mathbin{\overset{\mathbf L}{\otimes}}\mathbf Rf_*(A_X)
 \longrightarrow\mathbf Rf_*(A_X),\tag{2.1.1}
\]
它满足通常的\Term{associativity}{结合性}、\Term{commutativity}{交换性}，因而赋予 $\mathbf Rf_*(A_X)$ 一个“广义\Term{algebra}{代数}”的结构。为此，以
\[
 C^\bullet(X,A_X)
\]
\SourcePage{298}
表示 $A_X$ 的 Godement \Term{resolution}{分解}（SGA 4 第 XVII 讲\SGAsourcegap{原书此处在 XVII 后留白，未给出节号。}）。与 TF 6.6 中一样，定义\Term{resolution}{分解}的\Term{morphism}{态射}
\[
 C^\bullet(X,A_X)\otimes C^\bullet(X,A_X)\longrightarrow C^\bullet(X,A_X\otimes A_X).
\]
将它与从 $A_X$ 的\Term{multiplication}{乘法}得到的\Term{morphism}{态射}
\[
 C^\bullet(X,A_X\otimes A_X)\longrightarrow C^\bullet(X,A_X)
\]
复合，最终得到\Term{morphism}{态射}
\[
 C^\bullet(X,A_X)\otimes C^\bullet(X,A_X)\longrightarrow C^\bullet(X,A_X).\tag{2.1.2}
\]
应用\Term{functor}{函子} $f_*$，容易\Term{derive}{导出} $A_Y$-\Term{module}{模}\Term{complex}{复形}的\Term{morphism}{态射}
\[
 f_*C^\bullet(X,A_X)\otimes f_*C^\bullet(X,A_X)\longrightarrow f_*C^\bullet(X,A_X).
\]
再对左端各\Term{set-component}{分量}取\Term{flat}{平坦}\Term{resolution}{分解}，立即得到 (2.1.1)。

\textbf{2.2。} 设 $S$ 是\Term{scheme}{概形}，$E$ 是\Term{rank}{秩}为 $r+1$ 的\Term{locally-free}{局部自由} $\mathscr O_S$-\Term{module}{模}。以 $P=\mathbf P_S(E)\xrightarrow{p}S$ 表示相应的\Term{projective-bundle}{射影丛}。如所周知，由 Kummer \Term{exact-sequence}{正合列}，\Term{canonical}{典范} $\mathscr O_P$-\Term{module}{模} $\mathscr O_P(1)$ 确定一个元素
\[
 \xi\in H^2(P,\mu).
\]
以 $\eta$ 表示 $\xi$ 在\Term{canonical}{典范}\Term{morphism}{态射}
\[
 H^2(P,\mu)\longrightarrow H^0(S,R^2p_*(\mu))
\]
下的像，此\Term{morphism}{态射}来自 Leray \Term{spectral-sequence}{谱序列}
\[
 H^i(S,R^jp_*(\mu))\Longrightarrow H^{i+j}(P,\mu).
\]
类 $\xi$ 可视为 $D(P)$ 中的\Term{arrow}{箭头} $A_P\to\mu_P[2]$；\Term{transpose}{转置}后得到 $\mu_P^{\otimes-1}[-2]\to A_P$，从而由通常的\Term{adjunction}{伴随关系}得到\Term{arrow}{箭头}
\[
 c:\mu_S^{\otimes-1}[2]\longrightarrow\mathbf Rp_*(A_P),
\]
它在二维\Term{cohomology}{上同调}上与 $\eta$ 一致。

\SourcePage{299}
由\Term{tensor-product}{张量积}，借助 (2.2.1)，\Term{arrow}{箭头} $c$ 对每个整数 $i\ge1$ 确定\Term{arrow}{箭头}
\[
 c_i:\mu_S^{\otimes-i}[-2i]\longrightarrow\mathbf Rp_*(A_P),
\]
并约定以 $c_0$ 表示\Term{adjunction-morphism}{伴随态射} $A_S\to\mathbf Rp_*(A_P)$。这样就定义了 $D(S)$ 中的\Term{arrow}{箭头}
\[
 \gamma=\bigoplus_{i=0}^r c_i:
 \bigoplus_{i=0}^r\mu_S^{\otimes-i}[-2i]\longrightarrow\mathbf Rp_*(A_P).
\]

\begin{SGAresult}{\Term{theorem}{定理}}{2.2.1}
\Term{arrow}{箭头} $\gamma$ 是\Term{isomorphism}{同构}。
\end{SGAresult}
一旦定义了\Term{arrow}{箭头} $\gamma$，2.2.1 就\Term{equivalence}{等价}于下面的推论。

\begin{SGAresult}{命题}{2.2.2}
在上述假设下，有：

a) 对每个 $i\in\mathbf N$，$R^{2i+1}p_*(A_P)=0$。

b) 赋予 $\bigoplus_iR^{2i}p_*(\mu_P^{\otimes i})$ 从\Term{cup-product}{杯积}\Term{derive}{导出}的 $A_S$-\Term{algebraic-structure}{代数结构}后，由 $\eta$ 确定的 $A_S$-\Term{algebra}{代数}\Term{morphism}{态射}
\[
 A_S[T]\longrightarrow\bigoplus_iR^{2i}p_*(\mu_P^{\otimes i})
\]
的核是\Term{ideal}{理想} $(T^{r+1})$，因而它确定 $A_S$-\Term{algebra}{代数}\Term{isomorphism}{同构}
\[
 \theta:A_S[T]/(T^{r+1})\xrightarrow{\sim}\bigoplus_iR^{2i}p_*(\mu_P^{\otimes i}).
\]

c) 将 $\theta$ 与 $A_S[T]/(T^{r+1})$ 上的乘 $T^r$ 映射在 $A_S$ 上的\Term{restriction}{限制}复合，得到的\Term{morphism}{态射}
\[
 A_S\longrightarrow R^{2r}p_*(\mu_P^{\otimes r})
\]
是\Term{trace-morphism}{迹态射}（SGA 4 第 XVIII 讲）
\[
 R^{2r}p_*(\mu_P^{\otimes r})\longrightarrow A_S
\]
的逆；换言之，$\operatorname{Tr}(\eta^r)=1$。
\end{SGAresult}

\SourcePage{300}
\begin{proof}
这些数据在\Term{base-change}{基变换}下\Term{stable}{稳定}。由\Term{proper-base-change}{固有基变换定理}（SGA 4 第 XII 讲 5.1）以及\Term{trace-morphism}{迹态射}与\Term{base-change}{基变换}的相容性（SGA 4 第 XVIII 讲），可归结为在 $p$ 的\Term{fiber}{纤维}上验证相应断言。换言之，可假设 $S$ 为\Term{separably-closed}{可分闭}域 $k$ 的谱，$p$ 为\Term{canonical}{典范}\Term{morphism}{态射} $\mathbf P_k^r\to k$。此时证明要用下述\Term{lemma}{引理}。
\end{proof}

\begin{SGAresult}{\Term{lemma}{引理}}{2.2.3}
设 $k$ 是\Term{separably-closed}{可分闭}域，$p:\mathbf P_k^r\to k$ 为\Term{canonical-projection}{典范投影}。则：

a) 对每个 $i$，$R^{2i+1}p_*(A_P)=0$。

b) 当 $i\in[0,r]$ 时，$R^{2i}p_*(A_P)\simeq A_k$；当 $i>r$ 时，它为零。
\end{SGAresult}
\begin{proof}
\Term{lemma}{引理}在 $r=0$ 时显然成立，对 $r$ 归纳即可证明。假设它对 $r-1$ 成立，来证明 $r$ 的情形。\Term{projective-space}{射影空间} $\mathbf P_k^{r-1}$ 可作为\Term{closed-subscheme}{闭子概形}\Term{embedding}{嵌入} $\mathbf P_k^r$，其补\Term{open-subscheme}{开子概形}为 $\mathbf E_k^r$，由此得到\Term{compact-support-cohomology}{固有支集上同调}的\Term{exact-sequence}{正合列}
\[
 \cdots\longrightarrow H_!^i(\mathbf E_k^r)\longrightarrow H^i(\mathbf P_k^r)
 \longrightarrow H^i(\mathbf P_k^{r-1})\longrightarrow H_!^{i+1}(\mathbf E_k^r)\longrightarrow\cdots.
\]
由 $\mathbf E_k^r$ 的\Term{compact-support-cohomology}{固有支集上同调}计算（1.1），当 $i\notin\{2r,2r-1\}$ 时，\Term{inverse-image}{逆像}
\[
 H^i(\mathbf P_k^r)\longrightarrow H^i(\mathbf P_k^{r-1})
\]
是\Term{isomorphism}{同构}，故这一情形的断言成立。当 $i=2r-1$ 时，同样看到 $H^{2r-1}(\mathbf P_k^r)$ 包含于 $H^{2r-1}(\mathbf P_k^{r-1})$，后者由归纳假设为零。最后，当 $i=2r$ 时，由
\[
 H^{2r-1}(\mathbf P_k^{r-1})=H^{2r}(\mathbf P_k^{r-1})=0
\]
得到 $H_!^{2r}(\mathbf E_k^r)\xrightarrow{\sim}H^{2r}(\mathbf P_k^r)$，由 1.1 即得断言。
\end{proof}

\begin{proof}[续证]
\Term{lemma}{引理}\Term{pushout}{推出} 2.2.1 的 a) 部分。现在证明 $\theta$ 是\Term{isomorphism}{同构}。\Term{morphism}{态射} $\theta$ \Term{solvable}{可解}释为\Term{ring}{环}\Term{morphism}{态射}
\[
 A[T]/(T^{r+1})\longrightarrow\bigoplus_iH^{2i}(P,\mu^{\otimes i}).
\]
\SourcePage{301}
因此，由 2.2.3，只要证明对每个 $i\in[0,r]$，$\xi^i$ 是 $A$-\Term{module}{模} $H^{2i}(P,\mu^{\otimes i})$ 的自由元素即可。只须对 $\xi^r$ 验证这一点，而它立即由 2.2.2 c) \Term{pushout}{推出}。

所以来证明 c)。实际上，这就是当前情形中\Term{trace-morphism}{迹态射}的定义；无论如何，它表达了\Term{trace-morphism}{迹态射}与下图中\Term{fundamental-class}{基本类}的相容性：
\[
\begin{tikzcd}
 x_0 \arrow[r,hook,"\mathrm{inj.}"] \arrow[dr] & \mathbf P_k^r \arrow[d,"p"]\\
 & k
\end{tikzcd}
\]
其中 $\mathrm{inj.}:x_0\to\mathbf P_k^r$ 为 $\mathbf P_k^r$ 的一个\Term{closed-point}{闭点}的\Term{embedding}{嵌入}。
\end{proof}

\begin{SGAresult}{推论}{2.2.4}
设 $L$ 是一个 $A_S$-\Term{module}{模}\Term{complex}{复形}。\Term{arrow}{箭头}
\[
 \mathbf R\Gamma(S,\gamma\mathbin{\overset{\mathbf L}{\otimes}}\mathrm{id}(L)):
 \bigoplus_{i=0}^r\mathbf R\Gamma(S,L\otimes\mu_S^{\otimes-i})[-2i]
 \longrightarrow\mathbf R\Gamma(P,p^*(L))
\]
是\Term{isomorphism}{同构}。
\end{SGAresult}
事实上，只须对\Term{isomorphism}{同构} $\gamma\mathbin{\overset{\mathbf L}{\otimes}}_{A}\mathrm{id}_L$ 应用\Term{functor}{函子} $\mathbf R\Gamma_S$ 即可。

\textbf{2.2.5。} 为陈述下一个推论，对每个\Term{scheme}{概形} $X$，置
\[
 A^\bullet(X)=\bigoplus_iH^{2i}(X,\mu_X^{\otimes i}),
\]
以下始终采用此记号。

\begin{SGAresult}{推论}{2.2.6}
\Term{inverse-image}{逆像}\Term{homomorphism}{同态}
\[
 p^*:A^\bullet(S)=\bigoplus_iH^{2i}(S,\mu_S^{\otimes i})
 \longrightarrow\bigoplus_iH^{2i}(P,\mu_P^{\otimes i})=A^\bullet(P)
\]
是保持\Term{identity-element}{单位元}的\Term{ring}{环}的\Term{injective}{单射}\Term{morphism}{态射}，元素 $1,\xi,\ldots,\xi^r$ 构成 $A^\bullet(S)$-\Term{module}{模} $A^\bullet(P)$ 的一组基。特别地，$A^\bullet(P)$ 是一个\Term{free-module}{自由模}，其\Term{rank}{秩}与 $E$ 相同。
\end{SGAresult}

\SourcePage{302}
\SGAsection{3}{陈类}
设 $S$ 是\Term{scheme}{概形}，$E$ 是\Term{rank}{秩}恒为 $r$ 的\Term{locally-free}{局部自由} $\mathscr O_S$-\Term{module}{模}。以 $E^\vee$ 表示 $E$ 的\Term{dual}{对偶}，置 $P=\mathbf P_S(E^\vee)$。我们已经看到（2.2.6），若以 $\xi$ 表示 $\mathscr O_P(1)$ 在 $H^2(P,\mu)$ 中的类，则 $A^\bullet(S)$-\Term{module}{模} $A^\bullet(P)$ 为\Term{free-module}{自由模}，以 $1,\xi,\ldots,\xi^{r-1}$ 为基。特别地，由\Term{homogeneous}{齐次}性立即可知，恰有一个形如下式的关系：
\[
 \xi^r+c_1\xi^{r-1}+\cdots+c_i\xi^{r-i}+\cdots+c_r=0,\tag{3.1}
\]
其中对每个 $i\in[0,r]$，$c_i\in H^{2i}(S,\mu^{\otimes i})$。

现在设 $E$ 是任意\Term{locally-free}{局部自由} $\mathscr O_S$-\Term{module}{模}。对每个整数 $r$，$S$ 中 $E$ 的\Term{rank}{秩}为 $r$ 的点组成的\Term{set}{集合} $S_r$ 为\Term{open-subset}{开集}。因此得到将 $S$ 分成两两不交的\Term{open-subset}{开集}的分拆
\[
 S=\coprod_r S_r.
\]
特别地，对每个 $A_S$-\Term{module}{模} $M$ 及每个整数 $i$，各\Term{inverse-image}{逆像}之积给出的\Term{morphism}{态射}
\[
 H^i(S,M)\longrightarrow\prod_rH^i(S_r,M|_{S_r})
\]
为\Term{bijection}{双射}。

以上讨论说明，下述定义是合理的。

\begin{SGAresult}{定义}{3.2}
设 $S$ 是\Term{scheme}{概形}，$E$ 是\Term{locally-free}{局部自由} $\mathscr O_S$-\Term{module}{模}。对每个整数 $i\ge0$，$E$ 的第 $i$ 个\Term{chern-class}{陈类}记为 $c_i(E)$，它是 $H^{2i}(S,\mu^{\otimes i})$ 中的元素，其在每个整数 $r\ge0$ 所对应的\Term{open-subset}{开集} $S_r$ 上的\Term{restriction}{限制}为关系 (3.1) 所定义的元素 $c_i$；这里 $S_r$ 由 $E$ 的\Term{rank}{秩}为 $r$ 的点组成。若 $E$ 的\Term{rank}{秩}\Term{bounded}{有界}，则 $E$ 的\Term{total-chern-class}{总陈类}是 $A^\bullet(S)$ 中的元素
\[
 c(E)=\sum_i c_i(E).
\]
一般情形下，这个名称表示
\SourcePage{303}
\[
 \widehat A^\bullet(S)=\prod_iA^i(S)
\]
中的相应元素。
\end{SGAresult}

\textbf{3.3。} 显然，无论 $S,E$ 为何，均有 $c_0(E)=1$；若 $E$ 的\Term{rank}{秩}不超过 $d$，则当 $i>d$ 时 $c_i(E)=0$。

\begin{SGAresult}{命题}{3.4}
上面定义的\Term{chern-class}{陈类}满足下列性质。

(i)（规范化。）设 $L$ 是\Term{invertible-module}{可逆} $\mathscr O_S$-\Term{module}{模}。以 $[L]$ 表示 $L$ 的类在\Term{boundary-morphism}{边界态射}
\[
 H^1(S,\mathbf G_m)\longrightarrow H^2(S,\mu)
\]
下的像，则 $c(L)=1+[L]$。

(ii)（在\Term{base-change}{基变换}下\Term{stable}{稳定}。）设 $f:S'\to S$ 是\Term{scheme}{概形}的\Term{morphism}{态射}，$E$ 是\Term{locally-free}{局部自由} $\mathscr O_S$-\Term{module}{模}。对每个整数 $i$，有
\[
 f^*(c_i(E))=c_i(f^*E).
\]

(iii)（可\Term{additive}{加性}。）设 $0\to E'\to E\to E''\to0$ 是\Term{locally-free}{局部自由} $\mathscr O_S$-\Term{module}{模}的\Term{exact-sequence}{正合列}。对每个整数 $n$，有关系
\[
 c_n(E)=\sum_{i+j=n}c_i(E')c_j(E'').
\]
\end{SGAresult}
在开始证明之前，先指出下述推论。

\begin{SGAresult}{推论}{3.5}
对\Term{locally-free}{局部自由} $\mathscr O_S$-\Term{module}{模}的每个\Term{exact-sequence}{正合列} $0\to E'\to E\to E''\to0$，在 $\widehat A^\bullet(S)$ 中有
\[
 c(E)=c(E')c(E'').
\]
\end{SGAresult}

\textbf{3.4 的证明。}（依 A. Grothendieck：《\Term{chern-class}{陈类}理论》（La théorie des classes de Chern），\textit{Bull. S.M.F.}，1958，第 137--154 页。）

\SourcePage{304}
(i) 设 $E$ 是\Term{locally-free}{局部自由} $\mathscr O_S$-\Term{module}{模}；置 $P=\mathbf P_S(E^\vee)$，以 $p:P\to S$ 表示\Term{canonical-projection}{典范投影}。回忆有一个\Term{canonical}{典范}\Term{surjective}{满射}\Term{homomorphism}{同态}
\[
 p^*(E^\vee)\longrightarrow\mathscr O_P(1).
\]
特别地，当 $E$ \Term{invertible-module}{可逆}时，为与 (i) 的记号一致，置 $E=L$；此时 $P=S$，且有\Term{isomorphism}{同构}
\[
 L^\vee\xrightarrow{\sim}\mathscr O_P(1).
\]
所以 $\xi+[L]=0$，断言遂由\Term{chern-class}{陈类}的定义得出。

(ii) 直接验证即可，留给读者。

(iii) 将 $S$ 分成一些\Term{open-subset}{开集}，使所讨论的\Term{module}{模}在其上具有恒定\Term{rank}{秩}，便可假设 $E,E',E''$ 的\Term{rank}{秩}分别恒为 $r,r',r''$。暂且假定下面的\Term{lemma}{引理}成立，来说明它怎样\Term{pushout}{推出}所需断言。

\begin{SGAresult}{\Term{lemma}{引理}}{3.6}
设 $E$ 是一个\Term{locally-free}{局部自由} $\mathscr O_S$-\Term{module}{模}，带有\Term{finite}{有限}\Term{filtration}{滤过}
\[
 0=E_0\subset E_1\subset\cdots\subset E_i\subset\cdots\subset E_r=E,
\]
相邻的商 $L_i=E_i/E_{i-1}$（$1\le i\le r$）均\Term{invertible-module}{可逆}。则有
\[
 c(E)=(1+[L_1])(1+[L_2])\cdots(1+[L_r]).
\]
\end{SGAresult}
设 $h:Z\to S$ 是 $E'$ 与 $E''$ 的\Term{flag-scheme}{旗概形}的\Term{fiber-product}{纤维积}。由\Term{flag-scheme}{旗概形}的定义，$h^*(E')$ 和 $h^*(E'')$ 显然分别具有\Term{finite}{有限}\Term{filtration}{滤过}，其商分别为\Term{invertible-module}{可逆}\Term{module}{模} $L'_i$（$1\le i\le r'$）和 $L''_j$（$1\le j\le r''$）。于是 $h^*(E)$ 有一个\Term{filtration}{滤过}，其相邻的商是这些 $L'_i,L''_j$。应用\Term{lemma}{引理} 3.6 得到
\[
 c(h^*(E))=\prod_{1\le i\le r'}c(L'_i)\prod_{1\le j\le r''}c(L''_j)
 =c(h^*(E'))c(h^*(E'')),
\]
\SourcePage{305}
再由 (ii)，即得
\[
 h^*(c(E))=h^*(c(E')c(E'')).\tag{3.5.1}
\]
而\Term{morphism}{态射} $h$ 是一些形如
\[
 \mathbf P_T(F)\longrightarrow T\qquad(F\text{ 为局部自由的 }\mathscr O_T\text{-模})
\]
的\Term{morphism}{态射}的复合，所以（2.2.5）\Term{inverse-image}{逆像}\Term{homomorphism}{同态} $h^*:A^\bullet(S)\to A^\bullet(Z)$ 是\Term{injective}{单射}。由 (3.5.1) 得到所需断言。

现在证明\Term{lemma}{引理} 3.6。由上述论证，只须验证相应关系
\[
 c(p^*E)=\prod_i c(p^*(L_i)),
\]
其中仍以 $p:P=\mathbf P_S(E^\vee)\to S$ 表示\Term{canonical-projection}{典范投影}。对每个 $i\in[1,r]$，\Term{monomorphism}{单态射} $E_i\to E$ 经\Term{duality}{对偶}给出\Term{epimorphism}{满态射} $E^\vee\to E_i^\vee$；置 $Y_i=\mathbf P_S(E_i^\vee)$，得到 $P$ 的一列\Term{closed-subscheme}{闭子概形}：
\[
 S\simeq Y_1\subset Y_2\subset\cdots\subset Y_i\subset\cdots\subset Y_r=P.
\]
对每对满足 $1\le i\le j\le r$ 的整数 $(i,j)$，以 $\alpha_{ji}:Y_i\to Y_j$ 表示\Term{canonical-immersion}{典范}\Term{closed-immersion}{闭浸入}，并置 $\alpha_i=\alpha_{r,i}$。对每个整数 $i\in[1,r-1]$，由下面显然的\Term{lemma}{引理} 3.6.1，$Y_i$ 在 $Y_{i+1}$ 中由一个\Term{invertible-ideal}{可逆理想}定义，此\Term{ideal}{理想}\Term{isomorphism}{同构}于
\[
 \alpha_{i+1}^*(p^*(L_{i+1}^\vee)\otimes\mathscr O_P(-1)).
\]

\begin{SGAresult}{\Term{lemma}{引理}}{3.6.1}
设 $0\to L\to E\to F\to0$ 是\Term{locally-free}{局部自由} $\mathscr O_S$-\Term{module}{模}的\Term{exact-sequence}{正合列}。则
\[
 0\longrightarrow L\otimes_S\operatorname S(E)\xrightarrow{u}\operatorname S(E)
 \xrightarrow{v}\operatorname S(F)\longrightarrow0
\]
是\Term{graded-module}{分次} $\mathscr O_S$-\Term{module}{模}的\Term{exact-sequence}{正合列}，其中 $v$ 是显然的\Term{morphism}{态射}，$u$ 是与 $L$ 在 $E$ 中的包含对应的\Term{degree}{次数}为一的\Term{morphism}{态射}。
\end{SGAresult}

\SourcePage{306}
另一方面，$Y_1$ 可与 $S$ 认同，且有\Term{isomorphism}{同构}
\[
 \alpha_1^*(p^*(L_1^\vee))\simeq\mathscr O_{Y_1}(1),
\]
即
\[
 \alpha_1^*(p^*(L_1^\vee)\otimes\mathscr O_P(-1))\simeq\mathscr O_{Y_1}.\tag{3.6.2}
\]
我们将利用下述 Gysin \Term{homomorphism}{同态}的两个性质（SGA 4 第 XVIII 讲）作出结论。这里的 Gysin \Term{homomorphism}{同态}对应于\Term{smooth}{光滑} $S$-\Term{scheme}{概形}之间的\Term{closed-immersion}{闭浸入} $u:Z\to X$，它由 $\mathscr O_X$ 的\Term{ideal}{理想} $J$ 定义。

它满足\Term{projection-formula}{投影公式}
\[
 u_*(x\cdot u^*(y))=u_*(x)\cdot y.\tag{3.6.3}
\]
若\Term{ideal}{理想} $J$ \Term{invertible-module}{可逆}，则有关系
\[
 u_*(1_Z)=-[J]=[J^\vee].\tag{3.6.4}
\]
事实上，由此对整数 $i$ 作递增归纳，便可证明
\[
 (\alpha_i)_*(\alpha_i)^*\left(\prod_{j\le i}[p^*(L_j^\vee)\otimes\mathscr O_P(-1)]\right)=0.\tag{3.6.5}
\]
当 $i=r$ 时，断言 (3.6.5) 为
\[
 \prod_i[p^*(L_i^\vee)\otimes\mathscr O_P(-1)]=0,
\]
亦即
\[
 \prod_i\bigl(\xi+c_1(p^*(L_i^\vee))\bigr)=0,
\]
从而立即得到 3.6。$i=1$ 时，断言 (3.6.5) 由 (3.6.2) 得出。假设它对 $i$ 成立，来证明 $i+1$ 的情形。为简化记号，置
\[
 N=\prod_{j\le i}[p^*(L_j^\vee)\otimes\mathscr O_P(-1)].
\]
对于 $i$ 的等式 (3.6.5) 又可写为
\[
 (\alpha_{i+1})_*(\alpha_{i+1,i})_*(\alpha_{i+1,i})^*(\alpha_{i+1})^*(N)=0.
\]
再对 $\alpha_{i+1,i}$ 使用\Term{projection-formula}{投影公式}，得到
\SourcePage{307}
\[
 (\alpha_{i+1})_*\left[(\alpha_{i+1})^*(N)\cdot
 (\alpha_{i+1})^*[p^*(L_{i+1}^\vee)\otimes\mathscr O_P(-1)]\right]=0,
\]
所需结论随即成立。

\begin{SGAresult}{推论}{3.7}
(i) 对每个\Term{locally-free}{局部自由} $\mathscr O_S$-\Term{module}{模} $E$ 及每个整数 $i$，有 $c_i(E^\vee)=(-1)^ic_i(E)$。

(ii) 若 $E$ 是\Term{rank}{秩}为 $r$ 的\Term{locally-free}{局部自由} $\mathscr O_S$-\Term{module}{模}，则
\[
 c_1(E)=c_1(\mathop{\bigwedge}\nolimits^r E).
\]

(iii) 更一般地，\Term{tensor-product}{张量积}或\Term{exterior-power}{外幂}的\Term{chern-class}{陈类}可用 SGA 6 第 V 讲 6.1 中给出的\Term{universal-polynomial}{泛多项式}来表示。
\end{SGAresult}
\begin{proof}
例如证明 (ii)，其余断言同理。利用 $E$ 的\Term{flag-scheme}{旗概形}，可归结为 $E$ 有一个商均为\Term{invertible-module}{可逆}\Term{module}{模} $L_i$ 的\Term{filtration}{滤过}的情形。此时断言由熟知的\Term{isomorphism}{同构}
\[
 \mathop{\bigwedge}\nolimits^r(E)\simeq\bigotimes_iL_i
\]
得出。
\end{proof}

\textbf{3.8。} 给定 $\operatorname{Spec}(\mathbf C)$ 上\Term{locally-finite-type}{局部有限型}的\Term{scheme}{概形} $X$，其\Term{rational-point}{有理点}集 $X(\mathbf C)$ 自然带有一个\Term{locally-compact}{局部紧}\Term{topological-space}{拓扑空间}结构。因此，由\Term{algebra}{代数}\Term{topology}{拓扑}有一套\Term{chern-class}{陈类}理论 [3]，它为 $X(\mathbf C)$ 上的每个\Term{continuous}{连续} $\mathbf C$-\Term{vector-bundle}{向量丛} $V$ 赋予类 $c_i(V)\in H^{2i}(X(\mathbf C),\mathbf Z)$。现在以 $\mathscr O_c$ 表示 $X(\mathbf C)$ 上取复数值的\Term{continuous}{连续}函数\Term{sheaf}{层}，设 $\mathscr E$ 为\Term{locally-free}{局部自由}、\Term{finite-type}{有限型}的 $\mathscr O_c$-\Term{module}{模}；定义\Term{chern-class}{陈类}
\[
 c_i(\mathscr E)\in H^{2i}(X(\mathbf C),\mathbf Z),
\]
即以 $\mathscr E$ 为\Term{continuous}{连续}\Term{section}{截面}\Term{sheaf}{层}的 $\mathbf C$-\Term{vector-bundle}{向量丛}的\Term{chern-class}{陈类}。最后利用显然的映射 $\mathbf Z\to\mathbf Z/\nu\mathbf Z$，得到\Term{chern-class}{陈类}；在不致混淆时仍记为
\SourcePage{308}
\[
 c_i(\mathscr E)\in H^{2i}(X(\mathbf C),\mathbf Z/\nu\mathbf Z).\tag{3.8.1}
\]
另一方面，赋予 $X_{\text{ét}}$ 其\Term{structure-sheaf}{结构层} $\mathscr O_a$、赋予 $X(\mathbf C)$ \Term{ring}{环}\Term{sheaf}{层} $\mathscr O_c$ 后，有一个\Term{ringed-topos}{环化拓扑斯}\Term{morphism}{态射}（SGA 4 第 XI 讲）
\[
 \varepsilon:X(\mathbf C)\longrightarrow X_{\text{ét}}.
\]
其结构\Term{ring}{环}\Term{morphism}{态射}
\[
 \varepsilon^*(\mathscr O_a)\longrightarrow\mathscr O_c\tag{3.8.2}
\]
来自“每个\Term{algebra}{代数}函数均\Term{continuous}{连续}”这一事实。\Term{morphism}{态射} $\varepsilon$ 诱导\Term{abelian-sheaf}{阿贝尔层}的\Term{morphism}{态射}
\[
 \varepsilon^*(\mu_{X_{\text{ét}}})\longrightarrow\mu_{X(\mathbf C)},\tag{3.8.3}
\]
其中字母 $\mu$ 在两种情形中都表示 $\nu$ 次\Term{root-of-unity}{单位根}\Term{sheaf}{层}。最后，我们借助\Term{isomorphism}{同构}
\[
 (\mathbf Z/\nu\mathbf Z)_{X(\mathbf C)}\xrightarrow{\sim}\mu_{X(\mathbf C)}\tag{3.8.4}
\]
将两者作\Term{canonical}{典范}认同，此\Term{isomorphism}{同构}将\Term{section}{截面} $1$ 映为\Term{section}{截面} $e^{2i\pi/\nu}$。

\begin{SGAresult}{命题}{3.8.5}
设 $E$ 是\Term{locally-free}{局部自由} $\mathscr O_X$-\Term{module}{模}。对每个整数 $i\ge0$，以
\[
 \varepsilon^*:H^{2i}(X_{\text{ét}},\mu_X^{\otimes i})
 \longrightarrow H^{2i}(X(\mathbf C),\mathbf Z/\nu\mathbf Z)
\]
表示从 (3.8.3)、(3.8.4) 得到的\Term{inverse-image}{逆像}\Term{morphism}{态射}。则有等式
\[
 c_i(\varepsilon^*E)=\varepsilon^*c_i(E).
\]
\end{SGAresult}
\begin{proof}
只须对\Term{invertible-module}{可逆}的 $E$ 验证此式。此时，在通常的认同下，$\varepsilon^*(E)$ 是 $E$ 在\Term{canonical-map}{典范映射}
\[
 \varepsilon^*:H^1(X_{\text{ét}},\mathbf G_m)\longrightarrow H^1(X(\mathbf C),\mathscr O_c^*)\tag{3.8.6}
\]
\SourcePage{309}
下的像，所要验证的是下图交换：
\[
\begin{tikzcd}
 H^1(X,\mathbf G_m) \arrow[r,"\varepsilon^*"] \arrow[d,"\partial"'] & H^1(X(\mathbf C),\mathscr O_c^*) \arrow[d,"c_1"]\\
 H^2(X,\mu_X) \arrow[r,"\varepsilon^*"] & H^2(X(\mathbf C),\mathbf Z/\nu\mathbf Z).
\end{tikzcd}\tag{3.8.7}
\]
其中 $\partial$ 是由\Term{canonical}{典范}\Term{exact-sequence}{正合列}
\[
 0\longrightarrow\mu_X\longrightarrow\mathbf G_m\xrightarrow{\nu}\mathbf G_m\longrightarrow0\tag{3.8.8}
\]
得到的\Term{boundary-morphism}{边界态射}，$c_1$ 是 (3.8.1) 中的第一\Term{chern-class}{陈类}。事实上，利用“第一\Term{chern-class}{陈类}”\Term{morphism}{态射}
\[
 H^1(X(\mathbf C),\mathscr O_c^*)\longrightarrow H^2(X,\mathbf Z)
\]
是\Term{canonical}{典范}\Term{exact-sequence}{正合列}
\[
 0\longrightarrow\mathbf Z\xrightarrow{\ t\mapsto2i\pi t\ }\mathscr O_c
 \xrightarrow{\exp}\mathscr O_c^*\longrightarrow0
\]
的\Term{boundary-morphism}{边界态射}，容易看到 (3.8.7) 中的\Term{morphism}{态射} $c_1$ 正是\Term{exact-sequence}{正合列}
\[
 0\longrightarrow\mathbf Z/\nu\mathbf Z
 \xrightarrow{\ t\mapsto e^{2i\pi t/\nu}\ }\mathscr O_c^*
 \xrightarrow{\nu}\mathscr O_c^*\longrightarrow0
\]
所给出的\Term{boundary-morphism}{边界态射}。考虑到 (3.8.4)，这正是 (3.8.8) 的对应\Term{form}{形式}。断言随即成立。
\end{proof}

\textbf{3.9。} 设 $X$ 是\Term{scheme}{概形}，$E$ 是\Term{locally-free}{局部自由} $\mathscr O_X$-\Term{module}{模}。直接取\Term{limit}{极限}，即可对每个与 $X$ 的\Term{residue-characteristic}{剩余特征}互素的素数 $\ell$ 定义 $\ell$ 进\Term{chern-class}{陈类}
\[
 c_i(E)\in H^{2i}(X,\mathbf Z_\ell(i))
 \overset{\mathrm{def}}=\varprojlim_\nu H^{2i}(X,\mu_{\ell^\nu}^{\otimes i}).
\]
\SourcePage{310}
它们具有 3.4 和 3.7 中列出的通常性质，而且在 3.8 的情形中，显然与\Term{topology}{拓扑}方法定义的\Term{chern-class}{陈类}相容（相容的含义是显然的）。

再与 $\mathbf Q_\ell$ 作\Term{tensor-product}{张量积}，得到模去\Term{torsion}{挠元}的 $\ell$ 进\Term{chern-class}{陈类}
\[
 c_i(E)\in H^{2i}(X,\mathbf Q_\ell(i))
 =H^{2i}(X,\mathbf Z_\ell(i))\otimes_{\mathbf Z_\ell}\mathbf Q_\ell,
\]
它们显然满足相同的性质。

现在假设 $X$ 是域 $k$ 上的\Term{quasi-projective}{拟射影}\Term{smooth}{光滑}\Term{scheme}{概形}，在这里以 $C^\bullet(X)$ 表示其\Term{chow-ring}{Chow 环}（Chevalley 讨论班：《Chow 环及其应用》（Anneaux de Chow et applications），第 III 讲 4）。回忆我们在第 IV 讲中定义了一个\Term{graded-ring}{分次环}\Term{morphism}{态射}
\[
 \varphi_X:C^\bullet(X)\longrightarrow A^\bullet(X),
\]
它与\Term{direct-image}{正像}和\Term{inverse-image}{逆像}\Term{morphism}{态射}相容。另一方面，\Term{canonical}{典范}\Term{isomorphism}{同构}
\[
 \varepsilon_X:\operatorname{Pic}(X)\xrightarrow{\sim}C^1(X)
\]
将每个\Term{invertible-module}{可逆} $\mathscr O_X$-\Term{module}{模} $L$ 对应到 $L$ 的一个有理\Term{section}{截面}的\Term{divisor}{除子}的类，该\Term{section}{截面}在 $X$ 的任何\Term{connected-component}{连通分支}上都不恒为零。仿照上述 Grothendieck 的文章（其中假设 $k$ \Term{algebraically-closed}{代数闭}），这给出取值于 $C^\bullet(X)$ 的\Term{chern-class}{陈类}理论。复合\Term{morphism}{态射}
\[
 \varphi_X\circ\varepsilon_X:\operatorname{Pic}(X)\longrightarrow H^2(X,\mathbf Z_\ell(1))
\]
正是第一 $\ell$ 进\Term{chern-class}{陈类}。事实上，只须对 $\operatorname{Pic}(X)$ 中形如 $J^\vee$ 的元素验证，其中 $J$ 为 $\mathscr O_X$ 的一个\Term{invertible-ideal}{可逆理想}；此时由 (3.6.4) 即得结论。由\Term{chern-class}{陈类}理论的唯一性，遂知\Term{morphism}{态射} $\varphi_X$ 将\Term{chern-class}{陈类}映为\Term{chern-class}{陈类}。

\SourcePage{311}
\SGAsection{4}{自交公式及其应用}
设 $S$ 是\Term{scheme}{概形}，$X$ 是一个 $S$-\Term{scheme}{概形}，$Y$ 是 $X$ 的\Term{closed-subscheme}{闭子概形}，它在 $S$ 上\Term{smooth}{光滑}，且处处\Term{codimension}{余维数}为 $d$。还假设 $X$ 在 $Y$ 的某\Term{neighborhood}{邻域}内于 $S$ 上\Term{smooth}{光滑}。以 $i:Y\to X$ 表示\Term{canonical}{典范}\Term{closed-immersion}{闭浸入}，以 $J$ 表示定义 $Y$ 的 $\mathscr O_X$ 的\Term{ideal}{理想}。回忆在此情形下存在 Gysin \Term{isomorphism}{同构}（SGA $4\frac12$，Cycle）
\[
 i_*:H^{2p}(Y,\mu_Y^{\otimes p})\longrightarrow H_Y^{2p+2d}(X,\mu_X^{\otimes(p+d)}),
\]
即与\Term{fundamental-class}{基本类}的\Term{cup-product}{杯积}。在不致混淆时，对由忘掉\Term{support}{支集}而从上述\Term{morphism}{态射}得到的 Gysin \Term{morphism}{态射}
\[
 i_*:H^{2p}(Y,\mu_Y^{\otimes p})\longrightarrow H^{2p+2d}(X,\mu_X^{\otimes(p+d)})
\]
也采用同一记号。

\begin{SGAresult}{\Term{theorem}{定理}}{4.1}
在上述情形下，以 $N=J/J^2$ 表示 $i$ 的\Term{normal-sheaf}{法层}。则对每个 $y\in A^\bullet(Y)$，有等式
\[
 i^*i_*(y)=y\,c_d(N^\vee).
\]
\end{SGAresult}
在证明此\Term{theorem}{定理}之前，先指出几个变体，它们的证明完全相同。

a) 若 $\ell$ 是与 $S$ 的\Term{residue-characteristic}{剩余特征}互素的素数，则相应公式在 $H^*(Y,\mathbf Z_\ell(*))$ 中成立。

b) 可以加入\Term{support}{支集}。例如，若 $T$ 是 $Y$ 的\Term{closed-subset}{闭集}，$y\in H_T^{2p}(Y,\mu_Y^{\otimes p})$，则公式 4.1 仍在 $H_T^{2p+2d}(Y,\mu_Y^{\otimes(p+d)})$ 中成立。此时 $i_*$ 当然表示 Gysin \Term{morphism}{态射}
\[
 H_T^{2p}(Y,\mu_Y^{\otimes p})\longrightarrow H_T^{2p+2d}(X,\mu_X^{\otimes(p+d)}),
\]
\SourcePage{312}
它仍为\Term{isomorphism}{同构}（SGA $4\frac12$，Cycle）。

c) 最后，可以将以上两种情形结合起来……

4.1 的证明所用的基本工具是\Term{blow-up}{爆破}\Term{scheme}{概形}理论，因此先作一些回顾。以 $f:X'\to X$ 表示沿 $Y$ \Term{blow-up}{爆破}得到的 $X$-\Term{scheme}{概形}，考虑由 $f,i$ 构成的\Term{cartesian-diagram}{笛卡尔图}
\[
\begin{tikzcd}
 Y' \arrow[r,hook,"j"] \arrow[d,"g"'] & X' \arrow[d,"f"]\\
 Y \arrow[r,hook,"i"'] & X.
\end{tikzcd}
\]
回忆（参见 SGA 6 第 VII 讲 3.1，并沿用其中记号），$j$ 是\Term{codimension}{余维数}为 $1$ 的\Term{regular-immersion}{正则浸入}，由\Term{ideal}{理想}
\[
 J'\simeq\mathscr O_{X'}(1)
\]
定义，而且 $Y'\simeq\mathbf P_Y(N)$。另有 $\mathscr O_{Y'}$-\Term{module}{模}的\Term{exact-sequence}{正合列}
\[
 0\longrightarrow F\longrightarrow g^*(N)\longrightarrow\mathscr O_{Y'}(1)\longrightarrow0.
\]

\begin{SGAresult}{\Term{lemma}{引理}}{4.2}
对每个 $A_X$-\Term{module}{模} $M$，\Term{homomorphism}{同态}
\[
 f^*:H_Y^*(X,M)\longrightarrow H_{Y'}^*(X',f^*(M))
\]
与
\[
 f_*:H_{Y'}^*(X',f^*(M))\longrightarrow H_Y^*(X,M)
\]
满足关系 $f_*f^*=\mathrm{id}$。此外，$f_*$ 的\Term{degree}{次数}为 $0$。
\end{SGAresult}
\begin{proof}
因为\Term{morphism}{态射} $f$ 是\Term{relative-dimension}{相对维数}为零的\Term{relative-complete-intersection}{相对完全交}（SGA 6 第 VII 讲 1.8），最后一个断言由 Gysin \Term{homomorphism}{同态}的定义（SGA $4\frac12$，Cycle）得到。对于第一个断言，\Term{morphism}{态射} $f$ 的\Term{projection-formula}{投影公式}
\[
 f_*(x\cdot f^*(x'))=f_*(x)\cdot x',
\]
\SourcePage{313}
其中 $x\in H^*(X,A)$，$x'\in H_Y^*(X,M)$，使我们可归结为证明 $f_*(1)=1$。而只须在 $U=X\setminus Y$ 上验证相应断言即可；这显然成立，因为 $f$ 在 $U'=X'\setminus Y'$ 上诱导到 $U$ 的\Term{isomorphism}{同构}。
\end{proof}

\textbf{注记 4.3。} 当然，去掉\Term{support}{支集}后得到的与 4.2 相应的命题同样成立，证明也相同。

最后来证明 4.1。首先，由 SGA $4\frac12$，Cycle 中一个证明基本上是\Term{form}{形式}性的\Term{universal}{普遍}结果，有
\[
 i^*i_*(y)=y\,i^*i_*(1),
\]
所以只须证明
\[
 i^*i_*(1_Y)=c_d(N^\vee),
\]
其中 $i_*$ 在这里表示 Gysin \Term{isomorphism}{同构}
\[
 H^0(Y,A_Y)\longrightarrow H_Y^{2d}(X,\mu_X^{\otimes d}).
\]
因为 $f^*i_*(1_Y)\in H_{Y'}^{2d}(X',\mu_{X'}^{\otimes d})$，而 Gysin \Term{morphism}{态射}
\[
 j_*:H^{2d-2}(Y',\mu_{Y'}^{\otimes(d-1)})\longrightarrow H_{Y'}^{2d}(X',\mu_{X'}^{\otimes d})
\]
是\Term{isomorphism}{同构}，故存在 $u\in H^{2d-2}(Y',\mu_{Y'}^{\otimes(d-1)})$ 使
\[
 f^*i_*(1_Y)=j_*(u).\tag{4.1.1}
\]
对 (4.1.1) 两边应用 $j^*$，得到
\[
 j^*f^*i_*(1_Y)=j^*j_*(u).
\]
置 $\xi=c_1(\mathscr O_{Y'}(1))$，由 (3.6.4) 得到
\[
 g^*i^*i_*(1_Y)=-u\xi.\tag{4.1.2}
\]
元素 $u$ 有唯一的\Term{resolution}{分解}
\[
 \begin{split}
 u&=a_0+a_1\xi+\cdots+a_{d-1}\xi^{d-1},\\
 &\hspace{1em}a_i\in H^{2d-2-2i}(Y,\mu_Y^{\otimes(d-1-i)}).
 \end{split}\tag{4.1.3}
\]
考虑到 $N^\vee$ 的\Term{chern-class}{陈类}的定义关系
\SourcePage{314}
\[
 \xi^d+c_1(N^\vee)\xi^{d-1}+\cdots+c_{d-1}(N^\vee)\xi+c_d(N^\vee)=0,\tag{4.1.4}
\]
便得到等式
\[
 \begin{split}
 u\xi={}&-a_{d-1}c_d(N^\vee)
 +(a_0-a_{d-1}c_{d-1}(N^\vee))\xi+\cdots\\
 &+(a_{d-2}-a_{d-1}c_1(N^\vee))\xi^{d-1}.
 \end{split}
\]
将此式与 (4.1.2) 比较，并利用 $g^*$ 为\Term{injective}{单射}（2.2.6），得到
\[
 i^*i_*(1_Y)=a_{d-1}c_d(N^\vee).\tag{4.1.5}
\]
还须确定 $a_{d-1}$。首先，\Term{homomorphism}{同态} $g_*$ 将\Term{degree}{次数}降低 $2d-2$，故
\[
 g_*(\xi^i)=0\qquad(0\le i\le d-2).
\]
另一方面，由 2.2.2，$g_*(\xi^{d-1})=1$。对 (4.1.3) 两边应用 $g_*$，于是得到
\[
 g_*(u)=a_{d-1}.\tag{4.1.6}
\]
再对 (4.1.1) 两边应用 $f_*$，并考虑 4.2，得到
\[
 i_*(1_Y)=i_*g_*(u).\tag{4.1.7}
\]
这里的\Term{homomorphism}{同态} $i_*$ 为\Term{injective}{单射}，故
\[
 a_{d-1}=g_*(u)=1_Y,
\]
证明完成。

\textbf{注记 4.4。} a) 上述证明的关键在于\Term{purity-theorem}{纯性定理}。一旦在完全一般的情形下得到这个\Term{theorem}{定理}，4.1 就将对以下情形成立：$Y$ \Term{regular}{正则}，在 $X$ 中纯\Term{codimension}{余维数}为 $d$，且 $X$ \Term{excellent}{优良}、在 $Y$ 的某\Term{neighborhood}{邻域}内\Term{regular}{正则}。无论如何，由 SGA 4 第 XIX 讲，这个\Term{form}{形式}的结果在\Term{characteristic}{特征} $0$ 时已经成立。

b) 当 $S$ 是\Term{separably-closed}{可分闭}域的谱，而 $X,Y$
\SourcePage{315}
在 $S$ 上\Term{quasi-projective}{拟射影}且\Term{smooth}{光滑}时，$Y$ 的\Term{chow-ring}{Chow 环}中的相应命题已经由 Mumford 证明（参见 9.2）。

\begin{SGAresult}{推论}{4.5}
在 4.1 的假设下，设 $x,y\in A^\bullet(Y)$。则有
\[
 i_*(x)i_*(y)=i_*(xy\,c_d(N^\vee)).
\]
\end{SGAresult}
\begin{proof}
\Term{morphism}{态射} $i$ 的\Term{projection-formula}{投影公式}给出 $i_*(x)i_*(y)=i_*(x\,i^*i_*(y))$，再用 4.1 即得结论。
\end{proof}

\begin{SGAresult}{推论}{4.6}
设 $S$ 是\Term{scheme}{概形}，$E$ 是\Term{rank}{秩}恒为 $r$ 的\Term{locally-free}{局部自由} $\mathscr O_S$-\Term{module}{模}，$g:X=\mathbf V(E^\vee)\to S$ 是以 $E$ 为\Term{section}{截面}\Term{sheaf}{层}的\Term{vector-bundle}{向量丛}。则对 $X$ 在 $S$ 上的每个\Term{section}{截面} $f$，在 $H^{2r}(X,\mu_X^{\otimes r})$ 中有
\[
 f_*(1_S)=g^*(c_r(E)).
\]
\end{SGAresult}
\begin{proof}
由于 $g^*$ 是\Term{isomorphism}{同构}（1.2），且 $g\circ f=\mathrm{id}$，所以 $f^*$ 为\Term{bijection}{双射}。因而只须证明 $f^*f_*(1_S)=f^*g^*(c_r(E))$。以 $N$ 表示 $f$ 的\Term{normal-sheaf}{法层}，由 4.1，后一等式\Term{equivalence}{等价}于 $c_r(N^\vee)=c_r(E)$。事实上，甚至有\Term{isomorphism}{同构} $N^\vee\simeq E$：将 EGA IV 17.2.5 的\Term{exact-sequence}{正合列}用于 $S\xrightarrow{f}X\xrightarrow{g}S$，得到\Term{isomorphism}{同构} $N\simeq f^*\Omega^1_{X/S}$，另一方面 $\Omega^1_{X/S}\simeq g^*E^\vee$（EGA IV 16.4.8）。
\end{proof}

\begin{SGAresult}{推论}{4.7}
在 4.6 的假设下，再假设 $S$ 在某个\Term{separably-closed}{可分闭}域上\Term{smooth}{光滑}且\Term{quasi-projective}{拟射影}。以 $\sigma_X$ 表示 $X$ 的\Term{zero-section-cycle}{零截面圈}，以 $Y=f^{-1}(\sigma_X)$ 表示 $f$ 的\Term{zero-locus-cycle}{零点圈}，则在 $H^{2r}(S,\mu_S^{\otimes r})$ 中有
\[
 \operatorname{cl}_S(Y)=c_r(E).
\]
\end{SGAresult}

\SourcePage{316}
\begin{proof}
由第 IV 讲，$\operatorname{cl}_S(Y)=f^*\operatorname{cl}_X(\sigma_X)$。以 $s_0$ 表示 $X$ 的\Term{zero-section}{零截面}\Term{morphism}{态射}，有
\[
 \operatorname{cl}_X(\sigma_X)\overset{\mathrm{def}}=(s_0)_*(1_S)
 \overset{(4.6)}=g^*(c_r(E)).
\]
对两端应用\Term{homomorphism}{同态} $f^*$ 即得断言。
\end{proof}

\textbf{注记 4.8。} 设 $\ell$ 是不同于 $S$ 的\Term{residue-characteristic}{剩余特征}的素数。只要取\Term{inverse-limit}{逆极限}，就能由 4.6 和 4.7 得到 $\ell$ 进\Term{cohomology}{上同调}中的相应命题。

现在设 $\ell$ 是素数，$X$ 是\Term{separably-closed}{可分闭}域 $k$ 上的\Term{proper}{固有}\Term{scheme}{概形}，$k$ 的\Term{characteristic}{特征}与 $\ell$ 互素。回忆（第 IV 讲），$X$ 的 $\ell$ 进\Term{euler-poincare-characteristic}{Euler--Poincaré 特征}是整数
\[
 \operatorname{EP}_\ell(X)=\sum_i(-1)^i\dim H^i(X,\mathbf Q_\ell).
\]

\begin{SGAresult}{推论}{4.9}
设 $X$ 是\Term{separably-closed}{可分闭}域 $k$ 上\Term{dimension}{维数}为 $n$ 的\Term{smooth}{光滑}、\Term{proper}{固有}、\Term{connected}{连通}\Term{scheme}{概形}，且 $\operatorname{car}(k)\ne\ell$。以 $p:X\to k$ 表示\Term{canonical-projection}{典范投影}，则在 $\mathbf Q_\ell=H^0(k,\mathbf Q_\ell)$ 中有
\[
 \operatorname{EP}_\ell(X)=p_*c_n((\Omega_X^1)^\vee).
\]
\end{SGAresult}
\begin{proof}
设 $j:X\to X\times X$ 是 $X$ 的\Term{diagonal-immersion}{对角浸入}，$q:X\times X\to k$ 为\Term{canonical-projection}{典范投影}。对\Term{morphism}{态射} $j$ 应用 4.1 得到
\[
 c_n((\Omega_X^1)^\vee)=j^*j_*(1_X),
\]
故
\[
 p_*(c_n((\Omega_X^1)^\vee))=p_*j^*j_*(1_X)=q_*j_*j^*j_*(1_X).
\]
再利用\Term{morphism}{态射} $j$ 的\Term{projection-formula}{投影公式}，得等式
\[
 p_*(c_n((\Omega_X^1)^\vee))=q_*(j_*(1_X)\cup j_*(1_X)).
\]
其中由第 IV 讲，右端等于 $X$ 的 $\ell$ 进\Term{euler-poincare-characteristic}{Euler--Poincaré 特征}。
\end{proof}

\SourcePage{317}
\textbf{4.10。} 当 $X$ 在 $k$ 上\Term{projective}{射影}时，上面从第 IV 讲回忆的结果表明，$\operatorname{EP}_\ell(X)$ 是 $X\times X$ 的\Term{chow-ring}{Chow 环}中\Term{diagonal-cycle}{对角圈}的\Term{self-intersection}{自交}，因而不依赖于素数 $\ell\ne\operatorname{car}(k)$。于是以后可直接说 $X$ 的\Term{euler-poincare-characteristic}{Euler--Poincaré 特征} $\operatorname{EP}(X)$，不再提及整数 $\ell$。

\begin{SGAresult}{命题}{4.11}
设 $X$ 是\Term{separably-closed}{可分闭}域 $k$ 上\Term{dimension}{维数}为 $n$ 的\Term{smooth}{光滑}、\Term{projective}{射影}、\Term{connected}{连通}\Term{scheme}{概形}。则有
\[
 \operatorname{EP}(X)=\sum_{p,q}(-1)^{p+q}\dim H^p(X,\Omega_X^q).
\]
\end{SGAresult}
\begin{proof}
沿用 4.9 证明中的记号，以 $J$ 表示 $X$ 的\Term{diagonal-immersion}{对角浸入}的\Term{ideal}{理想}。如 4.10 所述，$\operatorname{EP}(X)$ 是 $X\times X$ 的\Term{chow-ring}{Chow 环}中\Term{diagonal-cycle}{对角圈} $\Delta$ 与自身的\Term{intersection-product}{交积} $\Delta\cdot\Delta$。为计算它，可以忽略\Term{torsion}{挠元}，因而可在 $\operatorname{Gr}_{\mathrm{top}}^\bullet(X)$ 中进行。这表明 $\operatorname{EP}(X)$ 是 $\mathscr O_{X\times X}/J$ 的类的平方在\Term{direct-image}{正像}\Term{morphism}{态射}
\[
 q_*:K(X\times X)\longrightarrow K(k)=\mathbf Z
\]
下的像。而由 SGA 6 第 VII 讲 2.5，这个平方等于
\[
 \sum_i(-1)^i\lambda^i([\Omega_X^1]).
\]
由 $q_*$ 的定义立即得到所需公式。
\end{proof}
